\section{Notation} \label{sec:notation}

% --------------------------------------------------------------------------

In this section, we prepare notation which we use throughout this paper.
The setting in Sec\-tion~\ref{sec:intro} is a particular case of the setting in this section.

We fix an integer $ h $ and a positive integer $ k $ such that $ \gcd(h, k) = 1 $.
Let $ M$, $N$, $a$, $c $ be positive integers and an integer $ b $ such that $ ac - MNb^2 = 1 $.
Let
\begin{gather*}
S = \pmat{Ma & MNb \\ MNb & Nc}
\end{gather*}
be a positive definite symmetric matrix.
The symbol $ \gamma = \trans{(\alpha, \beta)} $ denotes any element in $ \Q^2 $.
Here, $ t $ is the transpose of vectors.
Let $ Q(\alpha, \beta) := (\alpha, \beta) S \trans{(\alpha, \beta)} = Ma\alpha^2 + 2MNb \alpha \beta + Nc\beta^2 $ be the corresponding quadratic form.
Let $ \veps \colon (2S)^{-1}\big(\Z^2\big) \to (2S)^{-1}\big(\Z^2\big)/\Z^2 \to \bbC $ be a map.
For a complex number $ z $, we denote $ \bm{e}(z) := {\rm e}^{2\pi {\rm i} z} $.
Let $ \bbH $ be the upper half plane, that is, the set of complex numbers whose imaginary part are positive.
For a point $ \tau \in \bbH $, let {\samepage
\begin{gather*}
F_{Q, \veps}(\tau) := \sum_{0 \le \gamma \in (2S)^{-1}(\Z^2)} \veps(\gamma) q^{Q(\gamma)}
\end{gather*}
be a false theta function.}

We can rewrite a lattice $ S^{-1}\big(\Z^2\big) $ explicitly as below.
\begin{rem} \label{rem:S_inv_factor}
	Let
	\begin{gather*}
	A := \pmat{c & -Nb \\ -Mb & a} \in \SL_2(\Z).
	\end{gather*}
	Then we have
	\begin{gather*}
	SA = \pmat{M & 0 \\ 0 & N}
	\end{gather*}
	and $ S^{-1}(\Z^2) = \frac{1}{M}\Z \oplus \frac{1}{N}\Z $.
	Thus, we obtain $ \veps \colon \frac{1}{M}\Z/\Z \oplus \frac{1}{N}\Z/\Z \to \bbC $.
\end{rem}
\section{Vanishing results of Gauss sums} \label{sec:Gauss_sum}

In this section, we prove several claims which state vanishing results of weighted Gauss sums.
We~will use the results in this section to compute the WRT invariants and limit values of the homological blocks.
Let us keep the notation in Section \ref{sec:notation}, for example, $ h, k, Q(\gamma), \veps \colon (2S)^{-1}\big(\Z^2\big)/\allowbreak\Z^2 \to \bbC $.


\subsection{Main results in this section} %\label{subsec:Gauss_sum_main}

Throughout this section, we argue that Gauss sums vanish in the setting in Section~\ref{sec:notation} and assume the following assumption for a map $ \veps \colon (2S)^{-1}\big(\Z^2\big)/\Z^2 \to \bbC $.

\begin{assump}\qquad \label{assump:Gauss_sum}
	\begin{enumerate}
		\item[$(i)$] %\label{item:assump:Gauss_sum1}
		It holds
		\begin{gather*}
		\sum_{\gamma \in (2S)^{-1}(\Z^2)/\Z^2} \veps(\gamma) = 0.
		\end{gather*}
		\item[$(ii)$] %\label{item:assump:Gauss_sum2}
		For $ \gamma = \trans{(\alpha, \beta)} \in (2S)^{-1}\big(\Z^2\big) $ such that $ \veps(\gamma) \neq 0 $, the denominator of the irreducible fractions of $ \alpha $ and $ \beta $ are $ 2M $ and $ 2N $ respectively.
		\item[$(iii)$] %\label{item:assump:Gauss_sum3}
		For $ \gamma = \trans{(\alpha, \beta)} \in (2S)^{-1}\big(\Z^2\big) $ such that $ \veps(\gamma) \neq 0 $,
		\begin{gather*}
		M\alpha \pmod \Z, \qquad
		N\beta \pmod \Z, \qquad
		M\alpha^2 \pmod \Z,
\\
		N\beta^2 \pmod \Z, \qquad
		2MN\alpha \beta \pmod \Z
		\end{gather*}
		are independent of $ \gamma $.
	\end{enumerate}
\end{assump}

Here we remark that the map $ \veps \colon (2S)^{-1}\big(\Z^2\big)/\Z^2 \to \{ \pm1, 0 \} $ defined in Section \ref{sec:intro} satisfies the above assumption.

Our main results in this section are the following.

\begin{prop} \label{prop:Gauss_sum}
	Under Assumption $ \ref{assump:Gauss_sum} $, the followings are true:
	\begin{enumerate}
		\item[$(i)$] %\label{item:prop:Gauss_sum1} It holds
		It holds
\begin{gather*}
		\sum_{\gamma \in (2S)^{-1}(\Z^2)/k\Z^2} \veps(\gamma)
		\bm{e}\bigg( \frac{h}{k} Q(\gamma) \bigg)
		= 0.
		\end{gather*}
		\item[$(ii)$]% \label{item:prop:Gauss_sum2}
		We assume that there exist two maps $ \chi \colon \frac{1}{2M}\Z \to \frac{1}{2M}\Z/\Z \to \bbC $ and $ \psi \colon \frac{1}{2N}\Z \to \frac{1}{2N}\Z/\Z \to \bbC $ such that
		\begin{gather*}
		\sum_{\alpha \in \frac{1}{2M}\Z/\Z} \chi(\alpha) 		
		= \sum_{\beta \in \frac{1}{2N}\Z/\Z} \psi(\beta)
		= 0
		\end{gather*}
		and $ \veps\big(\trans{(\alpha, \beta)}\big) = \chi(\alpha) \psi(\beta) $ for any $ \alpha \in \frac{1}{2M}\Z/\Z $ and $ \beta \in \frac{1}{2N}\Z/\Z $.
		Then for any map $ C \colon \Q / k\Z \to \bbC $, we have
		\begin{gather*}
		\sum_{\gamma = \trans{(\alpha, \beta)} \in (2S)^{-1}(\Z^2)/k\Z^2} \veps(\gamma)
		\bm{e}\bigg(\frac{h}{k} Q(\gamma)\bigg) C(\alpha)
\\ \qquad
		{}=
		\sum_{\gamma = \trans{(\alpha, \beta)} \in (2S)^{-1}(\Z^2)/k\Z^2} \veps(\gamma)
		\bm{e}\bigg(\frac{h}{k} Q(\gamma)\bigg) C(\beta)
		= 0.
		\end{gather*}
		\item[$(iii)$] %\label{item:prop:Gauss_sum3}
		Under the assumption in $(ii)$, let $ B \colon \Q / k\Z \to \bbC $ be a map such that
		\begin{gather*}
		\sum_{\alpha \in \frac{1}{2M}\Z/\Z} \chi(\alpha) \widetilde{B}(\alpha) = 0,
		\end{gather*}
		where $ \widetilde{B}(\alpha) := \sum_{0 \le m < k} B(\alpha + m) $.
		Then it holds
		\begin{gather*}
			\sum_{\mu \in k\Z \oplus \Z/2kS(\Z^2)} \bm{e}\bigg(\frac{1}{4k} {\trans{\mu}} S^{-1} \mu\bigg)
			\sum_{\gamma = \trans{(\alpha, \beta)} \in (2S)^{-1}(\Z^2)/k\Z^2} \veps(\gamma)
			\bm{e}\bigg(\frac{1}{k} {\trans{\gamma}} \mu\bigg) B(\alpha) C(\beta)
\\ \qquad
		{}	=
			\sum_{\mu \in k\Z^2/2kS(\Z^2)} \bm{e}\bigg(\frac{1}{4k} {\trans{\mu}} S^{-1} \mu\bigg)\!\!
			\sum_{\gamma = \trans{(\alpha, \beta)} \in (2S)^{-1}(\Z^2)/k\Z^2}\!\! \veps(\gamma)
			\bm{e}\bigg(\frac{1}{k} {\trans{\gamma}} \mu\bigg) B(\alpha) C(\beta)
\\ \qquad
			{}= 0.
		\end{gather*}
		\item[$(iv)$] %\label{item:prop:Gauss_sum4}
		Under the assumption in $(ii)$, we have
		\begin{gather*}
			\sum_{\mu \in k\Z \oplus \Z/2kS(\Z^2)} \bm{e}\bigg(\frac{1}{4k} {\trans{\mu}} S^{-1} \mu\bigg)
			\sum_{\gamma = \trans{(\alpha, \beta)} \in (2S)^{-1}(\Z^2) \cap [0, k)^2} \veps(\gamma)
			\bm{e}\bigg(\frac{1}{k} {\trans{\gamma}} \mu\bigg) B_1(\alpha) B_1(\beta)
\\ \qquad
			{}=
			\sum_{\mu \in \Z \oplus k\Z/2kS(\Z^2)}\!\!\!\!\!\! \bm{e}\bigg(\frac{1}{4k} {\trans{\mu}} S^{-1} \mu\bigg)\!\!\!\!
			\sum_{\gamma = \trans{(\alpha, \beta)} \in (2S)^{-1}(\Z^2) \cap [0, k)^2} \!\!\!\!\!\veps(\gamma)
			\bm{e}\bigg(\frac{1}{k} {\trans{\gamma}} \mu\bigg) B_1(\alpha) B_1(\beta)
\\ \qquad
			{}=
			\sum_{\mu \in k\Z^2/2kS(\Z^2)}\!\!\!\! \bm{e}\bigg(\frac{1}{4k} {\trans{\mu}} S^{-1} \mu\bigg)\!\!\!\!
			\sum_{\gamma = \trans{(\alpha, \beta)} \in (2S)^{-1}(\Z^2) \cap [0, k)^2}\!\!\!\! \veps(\gamma)
			\bm{e}\bigg(\frac{1}{k} {\trans{\gamma}} \mu\bigg) B_1(\alpha) B_1(\beta)
\\ \qquad
			{}= 0,
		\end{gather*}
		where $ B_1(x) := x - 1/2 $ is the first Bernoulli polynomial.
		\item[$(v)$] %\label{item:prop:Gauss_sum5}
		Under the assumption in $(ii)$, we have
		\begin{gather*}
			\sum_{\mu \in k\Z \oplus \Z/2kS(\Z^2)} \bm{e}\bigg(\frac{1}{4k} {\trans{\mu}} S^{-1} \mu\bigg)
			\sum_{\gamma = \trans{(\alpha, \beta)} \in (2S)^{-1}(\Z^2) \cap [0, k)^2} \veps(\gamma)
			\bm{e}\bigg(\frac{1}{k} {\trans{\gamma}} \mu\bigg) \alpha \beta
\\ \qquad
			{}=
			\sum_{\mu \in \Z \oplus k\Z/2kS(\Z^2)} \bm{e}\bigg(\frac{1}{4k} {\trans{\mu}} S^{-1} \mu\bigg)
			\sum_{\gamma = \trans{(\alpha, \beta)} \in (2S)^{-1}(\Z^2) \cap [0, k)^2} \veps(\gamma)
			\bm{e}\bigg(\frac{1}{k} {\trans{\gamma}} \mu\bigg) \alpha \beta
\\ \qquad
			{}=
			\sum_{\mu \in k\Z^2/2kS(\Z^2)} \bm{e}\bigg(\frac{1}{4k} {\trans{\mu}} S^{-1} \mu\bigg)
			\sum_{\gamma = \trans{(\alpha, \beta)} \in (2S)^{-1}(\Z^2) \cap [0, k)^2} \veps(\gamma)
			\bm{e}\bigg(\frac{1}{k} {\trans{\gamma}} \mu\bigg) \alpha \beta
\\ \qquad
			{}= 0.
	\end{gather*}
	\end{enumerate}
\end{prop}

Our proof of Proposition \ref{prop:Gauss_sum} is based on the proof of \cite[Theorem 4.1]{BMM}.
However, our assumptions are slightly different from it in \cite[Theorem 4.1]{BMM}.

% --------------------------------------------------------------------------

\subsection[A proof of Proposition 4.2(i)]
{A proof of Proposition \ref{prop:Gauss_sum}$\boldsymbol{(i)}$}%\label{subsec:Gauss_sum1}

% --------------------------------------------------------------------------

Firstly, we prove Proposition \ref{prop:Gauss_sum}$(i)$.
To begin with, we need the result of the quadratic Gauss sum.
For integers $ a, b, c \in \Z $ with $ c > 0 $, we define the \textit{quadratic Gauss sum}
\begin{gather*}
G(a, b, c) := \sum_{n \in \Z/c\Z} \bm{e} \bigg(\frac{an^2 + bn}{c}\bigg).
\end{gather*}
The following result is elementary.

\begin{lem} \label{lem:Gauss_sum_fundamental}
	If $ \gcd(a, c) \nmid b $, then $ G(a, b, c) = 0 $.
\end{lem}

\begin{proof}
	Let $g := \gcd(a, c)$, $a' := a/g$, and $c' := c/g$.
	We have
	\begin{align*}
G(a, b, c)
&= \sum_{0 \le l < c'} \sum_{0 \le m < g} \bm{e} \bigg( \frac{ga'(l+c'm)^2 + b(l+c'm)}{gc'} \bigg)
\\
&= \sum_{0 \le l < c'} \bm{e} \bigg( \frac{ga' l^2 + bl}{gc'} \bigg) \sum_{0 \le m < g} \bm{e}
\bigg(\frac{bm}{g} \bigg).
	\end{align*}
	By assumption, the last sum vanishes.
\end{proof}

Proposition \ref{prop:Gauss_sum}$(i)$ follows from the following two lemmas.

\begin{lem} \label{lem:Gauss_sum_zero_general}
	If $ \gcd(M, k) > 1 $ or $ \gcd(N, k) > 1 $, then for $ \gamma = \trans{(\alpha, \beta)} \in \Q^2 $ we have
	\begin{gather*}
	\sum_{\mu \in \Z^2/k\Z^2} \bm{e}\bigg( \frac{h}{k} Q(\mu + \gamma) \bigg)
	= 0.
	\end{gather*}
\end{lem}

\begin{lem} \label{lem:Gauss_sum_indep_general}
	If $ \gcd(M, k) = \gcd(N, k) = 1 $, then
	\begin{gather*}
	\sum_{\mu \in \Z^2/k\Z^2} \bm{e}\bigg( \frac{h}{k} Q(\mu + \gamma) \bigg)
	\end{gather*}
	is independent of $ \gamma \in (2S)^{-1}\big(\Z^2\big) $ such that $ \veps(\gamma) \neq 0 $.
\end{lem}

\begin{proof}[Proof of Lemma $ \ref{lem:Gauss_sum_zero_general} $]
	We give a proof for the case when $ \gcd(M, k) > 1 $.
	Another case is similarly handled.
	We have
	\begin{align*}
		\sum_{\mu \in \Z^2/k\Z^2} \bm{e}\bigg( \frac{h}{k} Q(\mu + \gamma) \bigg)
		&= \sum_{\mu = \trans{(m, n)} \in \Z^2/k\Z^2} \bm{e}\bigg( \frac{h}{k} \big( Q(\mu) + 2\trans{\mu} S \gamma + Q(\gamma) \big) \bigg) \\
		&= \bm{e}\bigg(\frac{h}{k} Q(\gamma)\bigg)
		\sum_{n \in \Z/k\Z} \bm{e}\bigg( \frac{h}{k} \big( N c n^2 + 2Nn(Mb \alpha + c \beta) \big) \bigg) G_n,
	\end{align*}
	where $ r := 2M \alpha, s := 2N \beta \in \Z $, and $ G_n := G(hMa, h(ar + Mb(s + 2Nn)), k) $.
	Let $ g := \gcd(hMa, k) $.
	We have $ G_n = 0 $ if $ g \nmid h(ar + Mb(s + 2Nn)) $ by Lemma \ref{lem:Gauss_sum_fundamental}.
	Since $ g' := \gcd(M, k) $ divides $ g $, it suffices to show $ g \nmid h(ar + Mb(s + 2Nn)) $.
	
	We have $ \gcd(g', a) = \gcd(g', r) = \gcd(g', h) = 1 $ since $ ac - MN b^2 = 1 $ and $ \gcd(r, 2M) = \gcd(h, k) = 1 $ respectively.
	Thus, it holds $ \gcd(g, har) = 1 $.
	Therefore, we obtain
	\begin{gather*}
	h(ar + Mb(s + 2Nn)) \equiv ar \not\equiv 0 \pmod{g'}.\tag*{\qed}
\end{gather*}
\renewcommand{\qed}{}
\end{proof}

\begin{proof}[Proof of Lemma $ \ref{lem:Gauss_sum_indep_general} $]
	Since $ \gcd(M, k) = \gcd(N, k) = 1 $, there exist integers $ M^*, N^* \in \Z $ such that $ MM^* \equiv NN^* \equiv 1 \pmod k $.
	For $ \gamma \in (2S)^{-1}\big(\Z^2\big) $ such that $ \veps(\gamma) \neq 0 $,
	let $ \gamma^* := \trans{(MM^* \alpha, NN^* \beta)} \in \frac{1}{2} \Z^2 $.
	Since
	\begin{gather*}
	2 S (\gamma - \gamma^*)
	= 2 \pmat{Ma & MNb \\ MNb & Nc} \pmat{(1 - MM^*) \alpha \\ (1 - NN^*) \beta}
	\in k\Z^2,
	\end{gather*}
	we have
	\begin{gather*}
	Q(\mu + \gamma) - Q(\mu + \gamma^*)
	= Q(\gamma) - Q(\gamma^*) + 2 \trans{\mu} S (\gamma - \gamma^*)
	\equiv Q(\gamma) - Q(\gamma^*) \pmod{k\Z}.
	\end{gather*}
	Thus, we obtain
	\begin{align*}
		\sum_{\mu \in \Z^2/k\Z^2} \bm{e}\bigg( \frac{h}{k} Q(\mu + \gamma) \bigg)
		&= \sum_{\mu \in \Z^2/k\Z^2} \bm{e}\bigg( \frac{h}{k} ( Q(\mu + \gamma^*) + Q(\gamma) - Q(\gamma^*) ) \bigg)
\\
		&= \bm{e}\bigg( \frac{h}{k} ( Q(\gamma) - Q(\gamma^*)) \bigg)
		\sum_{\mu \in \Z^2/k\Z^2} \bm{e}\bigg( \frac{h}{k} Q(\mu) \bigg).
	\end{align*}
	By Assumption \ref{assump:Gauss_sum}$(iii)$, $ Q(\gamma) - Q(\gamma^*)\! \pmod{k\Z} $ is independent of $ \gamma $.
	Since $ \trans{(M\alpha, N\beta)} \!\!\pmod{\Z^2} $ is also independent of $ \gamma $ by Assumption \ref{assump:Gauss_sum}$(iii)$, there exists $ \delta_0 \in \frac{1}{2} \Z^2/\Z^2 $ which is independent of $ \gamma $ and $ \gamma^* \equiv \delta_0 \pmod{\Z^2} $.
	Thus,
	\begin{gather*}
	\sum_{\mu \in \Z^2/k\Z^2} \bm{e}\bigg( \frac{h}{k} Q(\mu + \gamma^*) \bigg)
	= \sum_{\mu \in \Z^2/k\Z^2} \bm{e}\bigg( \frac{h}{k} Q(\mu + \delta_0) \bigg)	
	\end{gather*}
	is also independent of $ \gamma $.
\end{proof}

\subsection[A proof of Proposition 4.2(ii)]{A proof of Proposition \ref{prop:Gauss_sum}$\boldsymbol{(ii)}$} %\label{subsec:Gauss_sum2}

Secondly, we prove Proposition \ref{prop:Gauss_sum}$(ii)$.

\begin{proof}[Proof of Proposition \ref{prop:Gauss_sum}$\boldsymbol{(ii)}$]
	We only prove the second equality.
	By assumption, we have
	\begin{gather*}
		\sum_{\gamma = \trans{(\alpha, \beta)} \in (2S)^{-1}(\Z^2)/k\Z^2} \veps(\gamma)
		\bm{e}\bigg(\frac{h}{k} Q(\gamma)\bigg) C(\beta)
\\ \qquad
		{}= \sum_{\beta \in \frac{1}{2N}\Z/k\Z} \psi(\beta)
		\bm{e}\bigg( \frac{h}{k} c \beta^2 \bigg) C(\beta)
		\sum_{\alpha \in \frac{1}{2M}\Z/k\Z} \chi(\alpha)
		\bm{e}\bigg( \frac{h}{k} M Q_{2N\beta}^{} (\alpha) \bigg),
	\end{gather*}
	where $ Q_s(\alpha) := a \alpha^2 + bs \alpha $.
	Then the second sum in the right hand side vanishes by the following lemma.
\end{proof}

\begin{lem} \label{lem:Gauss_sum_Bernoulli_single}
	Under the same assumption as in Proposition $\ref{prop:Gauss_sum}(ii)$, for a quadratic form $ Q_0(\alpha)\allowbreak = a_0 \alpha^2 + b_0 \alpha $ such that $ a_0, b_0 \in \Z, \gcd(M, a_0) = 1 $ and a map $ \widetilde{B} \colon \frac{1}{2M}\Z/\Z \to \bbC $ such that
	\begin{gather*}
	\sum_{\alpha \in \frac{1}{2M}\Z/\Z} \chi(\alpha) \widetilde{B}(\alpha) = 0,
	\end{gather*}
	it holds
	\begin{gather*}
	\sum_{\alpha \in \frac{1}{2M}\Z/k\Z} \chi(\alpha) \bm{e}\bigg( \frac{h}{k} M Q_0 (\alpha) \bigg) \widetilde{B}(\alpha) = 0.
	\end{gather*}
\end{lem}

Lemma \ref{lem:Gauss_sum_Bernoulli_single} follows from the following two lemmas.

\begin{lem} \label{lem:Gauss_sum_zero_Bernoulli}
	If $ \gcd(M, k) > 1 $, then for $ \alpha \in \frac{1}{2M} \Z $ such that $ \chi(\alpha) \neq 0 $, we have
	\begin{gather*}
	\sum_{m \in \Z/k\Z} \bm{e}\bigg( \frac{h}{k} M Q_0 (m + \alpha) \bigg)
	= 0.
	\end{gather*}
\end{lem}

\begin{lem} \label{lem:Gauss_sum_indep_Bernoulli}
	If $ \gcd(M, k) = 1 $, then
	\begin{gather*}
	\sum_{m \in \Z/k\Z} \bm{e}\bigg( \frac{h}{k} M Q_0 (m + \alpha ) \bigg)
	\end{gather*}
	is independent of $ \alpha \in \frac{1}{2M} \Z $ such that $ \chi(\alpha) \neq 0 $.
\end{lem}

\begin{proof}[Proof of Lemma $ \ref{lem:Gauss_sum_zero_Bernoulli} $]
	Let $ r := 2M\alpha $.
	We can write
	\begin{gather*}
	\sum_{m \in \Z/k\Z}\!\!\!\! \bm{e}\bigg( \frac{h}{k} M Q_0 \bigg(\!m \!+\! \frac{r}{2M} \bigg) \!\bigg)
\!= \bm{e}\bigg( \frac{h(a_0 r^2 \!+\! 2M b_0r)}{4Mk} \bigg)
\!\!\!	\sum_{m \in \Z/k\Z}\!\!\!\! \bm{e}\bigg( \frac{h(Ma_0 m^2 \!+\! (a_0r \!+\! Mb_0)m)}{k} \bigg).
	\end{gather*}
	Since $ \gcd(hMa, k) \mid \gcd(M, k) $ and $ \gcd(M, k) \nmid a_0 r + Mb_0 $, we obtain the claim.
\end{proof}

\begin{proof}[Proof of Lemma $ \ref{lem:Gauss_sum_indep_Bernoulli} $]
Since $ \gcd(M, k) = 1 $, there exists an integer $ M^* \in \Z $ such that $ MM^* \equiv 1 \pmod k $.
	For any integer $ m \in \Z $, we have
	\begin{gather*}
	Q_0 ( m + \alpha ) - Q_0 ( m + MM^* \alpha )
\\ \qquad
	{}	\equiv (2a_0 m + b)(1 - M M^*) \alpha + \big(1 - (MM^*)^2\big) a_0 \alpha^2 \pmod{k\Z}.
	\end{gather*}
	Thus, we obtain
	\begin{align*}
		\sum_{m \in \Z/k\Z}\!\! \bm{e}\bigg( \frac{h}{k} M Q_0 (m + \alpha ) \bigg)
		= {}
		&\bm{e}\bigg( \frac{h(1 - (MM^*)^2)}{k} a_0 M \alpha^2 + \frac{h(1 - M M^*)}{k} (2a_0 m + b) M
\alpha \bigg)
\\
		&\times \sum_{m \in \Z/k\Z} \bm{e}\bigg( \frac{h}{k} M Q_0 (m + MM^* \alpha ) \bigg).
	\end{align*}
	This sum is independent of $ \alpha $ by Assumption \ref{assump:Gauss_sum}$(iii)$.
\end{proof}

\subsection[A proof of Proposition 4.2(iii), (iv) and (v)]
{A proof of Proposition \ref{prop:Gauss_sum}$\boldsymbol{(iii)}$, $\boldsymbol{(iv)}$ and $\boldsymbol{(v)}$} %\label{subsec:Gauss_sum3}

Finally, we prove Proposition \ref{prop:Gauss_sum}$(iii)$, $(iv)$ and $(v)$.

\begin{proof}[Proof of Proposition \ref{prop:Gauss_sum}${\boldsymbol{(iii)}}$, $\boldsymbol{(iv)}$ and $\boldsymbol{(v)}$]
	Firstly, we prove Proposition~\ref{prop:Gauss_sum}$(iii)$.
	We on\-ly prove that the most left hand side vanishes.
	It can be written as
	\begin{gather*}
		\sum_{\beta \in \frac{1}{2N}\Z/k\Z}\!\!\!\! \psi(\beta) C(\beta)\!\!\!
		\sum_{n \in \Z/2kN\Z}\!\!\! \!\bm{e}\bigg( \frac{an^2}{4kN} + \frac{\beta n}{k} \bigg)
	\!\!\!\sum_{\alpha \in \frac{1}{2M}\Z/\Z}\!\!\! \!\chi(\alpha) \widetilde{B}(\alpha)\!\!\!
\sum_{m \in \Z/2M\Z}\!\!\!\!\bm{e}\bigg( \frac{kcm^2}{4M} - \frac{1}{2}bmn + \alpha m \bigg).
	\end{gather*}
	It suffices to show that the second double sum for $ \alpha $ and $ m $ vanishes.
	By Proposition \ref{prop:reciprocity}, it equals to
	{\samepage\begin{gather*}
\frac{2M{\rm i}}{\sqrt{kc}}
		\sum_{\alpha \in \frac{1}{2M}\Z/\Z} \chi(\alpha) \widetilde{B}(\alpha)
		\sum_{m \in \Z/kc\Z} \bm{e}\bigg({-}\frac{M}{kc} \bigg( m + \alpha - \frac{bn}{2} \bigg)^2 \bigg)
\\ \qquad
	{}=	\frac{2M{\rm i}}{\sqrt{kc}}
		\sum_{\alpha \in \frac{1}{2M}\Z/kc\Z} \chi\bigg(\alpha + \frac{Nbn}{2M} \bigg) \widetilde{B}\bigg(\alpha + \frac{Nbn}{2M} \bigg)
		\bm{e}\bigg({-}\frac{M}{kc} \alpha^2 \bigg).
	\end{gather*}
	This vanishes by Lemma \ref{lem:Gauss_sum_Bernoulli_single}.

}
	
	Secondly, we prove Proposition \ref{prop:Gauss_sum}$(iv)$.
	By the multiplication theorem of Bernoulli polynomials, we have
	\begin{gather*}
	B_1(\alpha) = \sum_{0 \le m < k} B\bigg( \frac{\alpha}{k} + m \bigg).
	\end{gather*}
	Since
	\begin{gather*}
	\sum_{\alpha \in \frac{1}{2M}\Z \cap [0, 1)} \chi(\alpha) B_1(\alpha)
	= \sum_{\alpha \in \frac{1}{2M}\Z \cap [0, 1)} \chi(\alpha) B_1(1-\alpha)
	= -\sum_{\alpha \in \frac{1}{2M}\Z \cap [0, 1)} \chi(\alpha) B_1(\alpha)
	\end{gather*}
	is equal to $ 0 $, we obtain the claim by Proposition \ref{prop:Gauss_sum}$(iii)$.
	
	Finally, we prove Proposition \ref{prop:Gauss_sum}$(v)$.
	By \ref{assump:Gauss_sum}$(i)$, we have
	\begin{align*}
		\sum_{\alpha \in \frac{1}{2M}\Z \cap [0, k)} \chi(\alpha) \alpha		&=
\sum_{\alpha \in \frac{1}{2M}\Z \cap [0, 1)} \chi(\alpha) \bigg( k\alpha + \frac{1}{2}k(k-1) \bigg)
		= 		k \sum_{\alpha \in \frac{1}{2M}\Z \cap [0, 1)} \chi(\alpha) \alpha
\\
		&= k \sum_{\alpha \in \frac{1}{2M}\Z \cap [0, 1)} \chi(\alpha) ( 1 - \alpha )
		= -k \sum_{\alpha \in \frac{1}{2M}\Z \cap [0, 1)} \chi(\alpha) \alpha.
	\end{align*}
	Since it vanishes, we obtain the claim by Proposition \ref{prop:Gauss_sum}$(iii)$.
\end{proof}

% --------------------------------------------------------------------------

\section{A limit value of a false theta function} \label{sec:false_theta_lim}

% --------------------------------------------------------------------------

In this section, we will calculate a radial limit of the false theta function $ F_{Q, \veps} (\tau) $.
We use notation in Section \ref{sec:notation} throughout this section.
A key fact is the following lemma which follows from the Euler--Maclaurin summation formula.

\begin{lem}[{\cite[Lemma 2.2]{BMM}}] \label{lem:Euler-Maclaurin}
	For two real numbers $ \alpha $ and $ \beta $ and $ C^\infty $-function $ f \colon \R^2 \to \bbC $ which has rapid decay, we have an asymptotic evaluation
	\begin{gather*}
	\sum_{m, n = 0}^{\infty} f(t(m+\alpha, n+\beta))
	\sim \sum_{j,l = -1}^{\infty} \frac{B_{j+1}(\alpha)}{(j+1)!} \frac{B_{l+1}(\beta)}{(l+1)!} f^{(j, l)}(0, 0) t^{j+l}
	\end{gather*}
	as $ t \to +0 $, where $ F(t) \sim G(t) $ means that $ F(t) = G(t) + O\big(t^N\big) $ for any positive integer $ N $,
	$ B_n(x) $ is the $ n $-th Bernoulli polynomial, and for integers $ j, l \ge 0 $ let
	\begin{gather*}
		f^{(-1, -1)}(0, 0) := \int_0^\infty \int_0^\infty f(x, y)\, {\rm d}x {\rm d}y, \\
		f^{(j, -1)}(0, 0) := -\int_0^\infty \frac{\partial^j f}{\partial x^j} (0, y)\, {\rm d}y, \\
		f^{(-1, l)}(0, 0) := -\int_0^\infty \frac{\partial^l f}{\partial y^l} (x, 0) \,{\rm d}x, \\
		f^{(j, l)}(0, 0) := \frac{\partial^{j+l} f}{\partial x^j \partial y^l} (0, 0).
	\end{gather*}
\end{lem}

We can calculate a radial limit of the false theta function $ F_{Q, \veps} (\tau) $ as follows.

\begin{lem} \label{lem:Phi_asymptotic}
	If
	\begin{gather*}
	\veps(\gamma) = \veps\big(\trans{(1, 1)} - \gamma\big) = \veps\big(\trans{(1 - \alpha, \beta)}\big) = \veps\big(\trans{(\alpha, 1 - \beta)}\big)
	\end{gather*}
	for any $ \gamma = \trans{(\alpha, \beta)} \in (2S)^{-1}\big(\Z^2\big) $, then we have an asymptotic evaluation
	\begin{equation} \label{eq:Phi_asymptotic}
		F_{Q, \veps} \bigg( \frac{h}{k} + \frac{t{\rm i}}{2\pi} \bigg)
		\sim \sum_{r = -1}^{\infty} a_{h, k}(r) t^r
	\end{equation}
	as $ t \to +0 $, where $ f(x, y) := {\rm e}^{-Q(x, y)} $ and
	\begin{align*}
		a_{h, k}(r) :={}& k^{2r}
\sum_{\gamma = \trans{(\alpha, \beta)} \in (2S)^{-1}(\Z^2) \cap [0, k)^2} \veps(\gamma) \bm{e}\bigg(\frac{h}{k} Q(\gamma)\bigg)
\\
		&\times\sum_{j+l = 2r, \, j, l \ge -1} \frac{1}{(j+1)! (l+1)!} B_{j+1}\bigg( \frac{\alpha}{k} \bigg) B_{l+1}\bigg( \frac{\beta}{k} \bigg) f^{(j, l)}(0, 0).
	\end{align*}
	
	Moreover, under Assumption $ \ref{assump:Gauss_sum} $ we have
	\begin{align*}
		a_{h, k}(r) ={}& k^{2r}
\sum_{\gamma = \trans{(\alpha, \beta)} \in (2S)^{-1}(\Z^2) \cap [0, k)^2} \veps(\gamma) \bm{e}\bigg(\frac{h}{k} Q(\gamma)\bigg) \\
		&{}\times\sum_{j+l = 2r, \, j, l \ge 0} \frac{1}{(j+1)! (l+1)!} B_{j+1}\bigg( \frac{\alpha}{k} \bigg) B_{l+1}\bigg( \frac{\beta}{k} \bigg) f^{(j, l)}(0, 0).
	\end{align*}
\end{lem}

\begin{proof}
	Our proof is based on the proof of \cite[Lemma 3.2]{BMM}.
	Since we can write
	\begin{gather*}
	F_{Q, \veps} \bigg( \frac{h}{k} + \frac{t{\rm i}}{2\pi} \bigg)
	=
	\sum_{\gamma = \trans{(\alpha, \beta)} \in (2S)^{-1}(\Z^2) \cap [0, k)^2} \veps(\gamma) \bm{e}\bigg(\frac{h}{k} Q(\gamma)\bigg)
	\sum_{\mu \in \Z_{\ge 0}^2} f(\sqrt{t}(\mu + \gamma)),
	\end{gather*}
	we have an asymptotic evaluation
	\begin{align*}
		F_{Q, \veps} \bigg( \frac{h}{k} + \frac{t{\rm i}}{2\pi} \bigg)
\sim{} &\sum_{j,l = -1}^{\infty} k^{j+l} t^{(j+l)/2}
\sum_{\gamma = \trans{(\alpha, \beta)} \in (2S)^{-1}(\Z^2) \cap [0, k)^2} \veps(\gamma) \bm{e}\bigg(\frac{h}{k} Q(\gamma)\bigg)
\\
		&\times\frac{1}{(j+1)! (l+1)!} B_{j+1}\bigg( \frac{\alpha}{k} \bigg) B_{l+1}\bigg( \frac{\beta}{k} \bigg) f^{(j, l)}(0, 0)
	\end{align*}
	by Lemma \ref{lem:Euler-Maclaurin}.
	By assumption, the right hand side yields
	\begin{gather*}
		\sum_{j,l = -1}^{\infty} k^{j+l} t^{(j+l)/2}
		\sum_{\gamma = \trans{(\alpha, \beta)} \in (2S)^{-1}(\Z^2) \cap [0, k)^2} \veps(\gamma) \bm{e}\bigg( \frac{h}{k} Q(k - \alpha, k - \beta) \bigg)
\\ \qquad
		{}\times\frac{1}{(j+1)! (l+1)!} B_{j+1}\bigg( \frac{k - \alpha}{k} \bigg) B_{l+1}\bigg( \frac{k - \beta}{k} \bigg) f^{(j, l)}(0, 0).
	\end{gather*}
	Since $ Q(k - \alpha, k - \beta) \equiv Q(\gamma) \pmod{k\Z} $ and $ B_n(1-x) = (-1)^n B_n(x) $, this becomes
	\begin{gather*}
		\sum_{j,l = -1}^{\infty} k^{j+l} t^{(j+l)/2} (-1)^{j+l+2}
		\sum_{\gamma = \trans{(\alpha, \beta)} \in (2S)^{-1}(\Z^2) \cap [0, k)^2} \veps(\gamma) \bm{e}\bigg(\frac{h}{k} Q(\gamma)\bigg)
\\ \qquad
		{}\times\frac{1}{(j+1)! (l+1)!} B_{j+1}\bigg( \frac{\alpha}{k} \bigg) B_{l+1}\bigg( \frac{\beta}{k} \bigg) f^{(j, l)}(0, 0).
	\end{gather*}
	Since it vanishes for the odd values of $ j+l $, we obtain the asymptotic evaluation (\ref{eq:Phi_asymptotic}).
	
	The latter claim follows from Proposition \ref{prop:Gauss_sum}$(i)$.
\end{proof}

Thus, we obtain the following proposition by Proposition \ref{prop:Gauss_sum}$(i)$ and $(ii)$.

\begin{prop} \label{prop:Phi_lim}
	If Assumption $ \ref{assump:Gauss_sum} $ and the assumption in Lemma $ \ref{lem:Phi_asymptotic} $ and Proposition $ \ref{prop:Gauss_sum}(ii)$ hold, then we have\vspace{-.5ex}
	\begin{gather*}
		\lim_{t \to 0} F_{Q, \veps} \bigg( \frac{h}{k} + ti \bigg)
		= \frac{1}{k^2}
		\sum_{\gamma = \trans{(\alpha, \beta)} \in (2S)^{-1}(\Z^2) \cap [0, k)^2} \veps(\gamma) \bm{e}\bigg(\frac{h}{k} Q(\gamma)\bigg)
		\alpha \beta.
	\end{gather*}
\end{prop}
